\documentclass[colorlinks,empty]{goose-article}

\usepackage{tcolorbox}

\title{Teaching}
\author{Tom de Geus}
\hypersetup{pdfauthor={T.W.J. de Geus}}

\begin{document}

\maketitle

\section*{Statistical Physics 2}

\begin{tcolorbox}
    \begin{itemize}
        \item Bachelor.
        \item Teaching assistant, EPFL, 02.2020 -- 07.2020.
        \\
        Responsibilities: exercise sessions and grading.
    \end{itemize}
\end{tcolorbox}

\noindent
Introduction to the theory of phase transitions.
The course covers the following topics:

\begin{enumerate}
    \item Phase transitions.
    \begin{itemize}
        \item Classification
        \item Liquid-gas transition
        \item Ferromagnetism
        \item Symmetry breaking
    \end{itemize}
    \item New techniques.
    \item Renormalisation Group (RG).
\end{enumerate}

\section*{Continuum Mechanics}

\begin{tcolorbox}
    \begin{itemize}
        \item Bachelor.
        \item Teaching assistant, EPFL, 09.2016 -- 02.2017.
        \\
        Responsibilities: exercise sessions and grading.
    \end{itemize}
\end{tcolorbox}

\noindent
Continuum mechanics introduces the concepts of stress and strain, conservation laws, equilibrium equations and constitutive laws.
The course covers the following topics:

\begin{itemize}
    \item Continuum approximations and conversation laws.
    \item Tensors.
    \item Stress and strain.
    \item Elasticity.
    \item Work and energy.
\end{itemize}

\clearpage
\section*{Finite Element Method}

\begin{tcolorbox}
    \begin{itemize}
        \item Bachelor.
        \item Teaching assistant, TU/e, 09.2010 -- 09.2011.
        \\
        Responsibilities: exercise sessions and grading.
        \item Co-developer of a tensor toolbox in Matlab for education purposes, TU/e, 09.2009 -- 09.2010.
    \end{itemize}
\end{tcolorbox}

\noindent
The finite element method is a numerical technique for finding (approximate) solutions to boundary value problems for partial differential equations.
The course covers the following topics:

\begin{itemize}
    \item Discretization and interpolation.
    \item Variational formulation.
    \item Numerical integration.
    \item Linear and non-linear problems.
    \item Implementation in Matlab.
\end{itemize}

\section*{Programming project}

\begin{tcolorbox}
    \begin{itemize}
        \item Bachelor.
        \item Lecturer, TU/e, 09.2011 -- 09.2015.
        \\
        Responsibilities: course development, independent teaching, and grading.
    \end{itemize}
\end{tcolorbox}

\begin{itemize}
    \item Gaining skills in anticipation, analysis and interpretation.
    \item Gaining skills in scientific presentation of experimental and numerical results and uncertainties.
    \item Gaining practical skills in the field of numerical modelling of solid-state mechanics.
    \item Graining practical programming skills, including pointers and data-structures.
\end{itemize}

\section*{Design Based Learning}

\begin{tcolorbox}
    \begin{itemize}
        \item Bachelor.
        \item Supervisor, TU/e, 09.2008 -- 09.2009.
        \\
        Responsibilities: project supervision and grading.
    \end{itemize}
\end{tcolorbox}

\noindent
In Design Based Learning, students apply knowledge from previous (theoretical) courses and learn to collaborate in multidisciplinary teams.
By doing so, students will not only develop (new) disciplinary knowledge and technical skills, but also improve on their professional and personal skills.
Bi-weekly team meetings are supervised by staff members that guide both the process and the content of the project.

\end{document}
